% SPDX-License-Identifier: Apache-2.0
% covers: ChanCdc
\datasheet{ChanCdc}{A channel from one clock to another, with Gray
coded pointers and two flops each way.}

\dsfacts{\code{lib/parts/src/cdc.rs}}{\code{txhdl\_parts::cdc::ChanCdc}}%
{\code{//docs:examples}}

\subsection*{Function}

Two clocks that are not the same clock have no fixed relation, so a
register written by one and read by the other is read at an arbitrary
moment in its life, including while it is changing. \code{ChanCdc}
carries words across that gap: a channel in on the writing clock, a
channel out on the reading clock, and a memory between them written
by one side and read by the other.

It is one unit with two processes, one per clock. What crosses
between them is only the two pointers, each through two flip-flops of
the clock that reads it.

\subsection*{Parameters}

\begin{center}\footnotesize
\begin{tabular}{@{}lp{0.68\textwidth}@{}}
\toprule
Parameter & Meaning \\
\midrule
\code{T} & What it carries: a \code{Transaction} and \code{Value} type. \\
\code{AW} & The width of the address, in bits. \\
\code{D} & How many words it holds, which must be \code{1 << AW}. \\
\code{PW} & The width of the pointers, which must be \code{AW + 1}. \\
\code{W} & The writing clock. \\
\code{R} & The reading clock. \\
\bottomrule
\end{tabular}
\end{center}

\code{D} and \code{PW} are stated rather than computed, because an
expression in a const parameter's position needs nightly Rust, which
is the same reason \code{Fifo} takes both \code{AW} and \code{D}. The
tables below use \code{ChanCdc<U<32>, 4, 16, 5, DefaultClock,
ClkPix>}, the crossing of \code{ex\_cdc}: sixteen words of thirty-two
bits.

\dstables{ChanCdc}

\subsection*{Behaviour}

\begin{itemize}
\item \code{wbin} and \code{rbin} are the two pointers in binary, each
on its own clock. \code{wgray} and \code{rgray} are the same pointers
Gray coded, and they are what the other side samples.
\item \code{rgray\_w1} and \code{rgray\_w2} are the reading pointer
caught on the writing clock, and \code{wgray\_r1} and \code{wgray\_r2}
the writing pointer caught on the reading clock. The second of each
pair is the one its side reads.
\item Empty is \code{rgray} equal to \code{wgray\_r2}. Full is
\code{wgray} equal to \code{rgray\_w2} in every bit but the top two.
\item A word is taken when one is offered and it is not full, and
lands in the memory at the writing pointer. A word is offered out
when it is not empty, read from the memory at the reading pointer.
\end{itemize}

The pointers are Gray coded because a binary pointer changes several
bits at once: \code{0111} to \code{1000} is four bits, and a sample
caught among them can read any of sixteen values, most of which the
pointer never held. One bit per step means a sample caught mid-change
reads the value before or the value after and nothing else. Both are
safe, and safe in a particular direction: the reader concludes the
queue is emptier than it is and the writer that it is fuller, so each
errs towards refusing work rather than corrupting it.

The pointers are one bit wider than the address so that full and empty
are distinguishable, which equal pointers of the same width are not.

\subsection*{Verification}

\begin{itemize}
\item \code{ex\_cdc} runs the writing side on a clock of period four
and the reading side on one of period six with a phase of two, so the
edges drift past each other and every relative position happens
somewhere in the run. The build checks the netlist against that run
under nvc and under Verilator:
\code{//docs:cdc\_sim\_cdc\_tb\_test} and
\code{//docs:cdc\_vsim\_test}.
\item Those checks reach inside. Every register is compared by
hierarchical name as well as every port, so both Gray pointers and all
four synchroniser flops are checked tick by tick, and a pointer that
advanced at the wrong moment or a full test off by the wrap bit is a
failure rather than a difference nobody looks at.
\end{itemize}

\subsection*{Limits}

\textbf{Metastability is not modelled and cannot be.} Both sides are
deterministic and the sampling clock reads a register that has
settled, so the two flops are never seen to do anything. What the
co-simulation demonstrates is agreement about the pointer protocol.
The second flop is an argument about silicon, and no simulation in
either language says anything about it. A passing test beside a
synchroniser invites the opposite conclusion, which is why this
paragraph is here rather than in a comment.

\code{D} must be \code{1 << AW} and \code{PW} must be \code{AW + 1},
and nothing checks either.

It does not replace \code{lib/board/chan\_cdc.v}, of which the
flagship holds five instances. That is a separate decision and a
separate change: the hand-written file is proven on hardware and this
one is proven in simulation, and those are not the same claim.
